% !TeX root = ../papers/nc-p3.tex
\section{有限両対称窓の非正則境界を与える二元十次形式と三割線}
\label{sec:binary-secant-obstructions}
$P_j=k[s,t]_j$ とし、$D\in P_{a+b}^*$ の積対合を
$H_{a,b}(D):P_b\to P_a^*$ と書く。以下の幾何的記述は代数閉包上で行う。
\begin{theorem}\label{thm:binary-secant-obstructions}
点グラフが有限な両対称Frobenius窓のうち非正則なものは、必要なら始点を一つ移し、
以下の構成で全て得られる。次数十の汎関数 $T$ は $G=H_{5,5}(T)$ が可逆なもの、
次数八の汎関数 $Y$ は $T\mathbin{\lrcorner}P_2$ の外にあるものを取る。
$U=P_2\oplus k$, $Q=P_3$ とおき、積を
\[
 B((p,a),q)=pq+aG^{-1}H_{5,3}(Y)q\in P_5
\]
とする。四窓は $w(u,q,u',q')=G(B(u,q),B(u',q'))$ である。
固定した二元座標と $[T]$ に対し、余分な文字の選び方を除いたパラメータは
$[\bar Y]\in\mathbb P(P_8^*/(T\mathbin{\lrcorner}P_2))=\mathbb P^5$ である。
ここで始点の移動は二つの向きを列挙するために用いる。
二元形式への商の印を忘れた同型類の重複は、このパラメータ表示では取り除いていない。

$S_j$ を有理正規八次曲線の長さ $j$ の割線多様体とする。
中心 $\Pi=\mathbb P(T\mathbin{\lrcorner}P_2)$ からの射影
$\pi:S_3\to\mathbb P^5$ は次数20の有限平坦被覆である。
点グラフが有限であるための必要十分条件は
$[\bar Y]\notin\pi(S_2)$ であり、そのときグラフは被覆のスキーム論的ファイバーである。
この開集合では定理\ref{thm:frobenius-tensor-reconstruction}の三次・四次交叉を調べれば
求める非正則Frobenius窓が正確に得られる。一般の $(T,[\bar Y])$ は両交叉検査を通り、
相異なる20点を持つ安定な非正則窓を与える。固定した $T$ での有限平坦被覆の
トレース判別式は $\mathbb P^5$ 上次数90の因子を定める（重複度を含む）。
\end{theorem}
\begin{proof}[二元商からの復元]
有限Frobenius窓では純関係がなく、両評価の階数は3以上である。
非正則なら定理\ref{thm:bisymmetric-two-length-test}により一方に二元形式への商がある。
関係六平面を商に写す核は、一方の文字が固定された純関係だけからなり零。
六次元どうしなので二元乗法の双対像に一致する。双対に移り、関係空間の双対を $P_5$ と同一視し、
指定された三文字の補空間を選ぶと積は $B((p,a),q)=pq+aLq$ である。
両対称性と隣接切断の階数6により $w$ は非退化な対称形式 $G$ によるこの積のGram形式である。
逆に、一方の $q,q'$ の交換不変性だけを課せば、Gram対称性から他方も交換不変となる。

二元単項式を $t$ の指数で番号付けする。三文字どうしの交換条件は
$G_{a+i,b+j}=G_{a+j,b+i}$ ($0\le a,b\le2$, $0\le i,j\le3$) であり、
対称行列 $G$ の解はちょうどHankel行列である。
$J=GL$ に対する条件は $J_{a+i,j}=J_{a+j,i}$ となり、解も全てHankelである。
実際どちらも座標どうしの同一視だけで、第一の21座標は11個、第二の24座標は9個の
反対角線ごとに連結する（添付検算では各等式の整数行列の階数10と15を求める）。
従って $G,J$ は一意に $T,Y$ を与える。
第四文字の平行移動と非零倍は $Y$ の三次元の極性空間による変更とスカラー倍となる。
極性空間に入れば第四文字が他の三つに従属する。外では四文字は独立であり、
$Q$ 側の独立性は二元乗法から従う。$B$ は全射で、両隣接階数は6である。
逆向きには二元商による指数成長が全てのパラメータで正則性を排除する。
\end{proof}
\begin{proof}[有限被覆と点グラフ]
$u=[p:a]$ での関係評価は $G$ を掛けると $H_{5,3}(T\mathbin{\lrcorner}p+aY)$ であり、
その核にある $[q]$ がグラフの他の座標になる。$a=0$ かつ $p\ne0$ なら
これは $q\mapsto G(pq)$ で階数4。従って中心は階数3以下の多様体と交わらない。
$([Z],[q])\in\mathbb P(P_8^*)\times\mathbb P(P_3)$ で $Z(qP_5)=0$ の入射を取る。
非零 $q$ による乗法は単射なので、これは $\mathbb P^3$ 上の階数3の束
$E=\ker(9\mathcal O\to6\mathcal O(1))$ の直線の射影化である。
相異なる三根を持つ一般の $q$ ではその平面は根の三つの評価点の張る平面であり、
一般の $Z$ は階数3でただ一つの核 $[q]$ を持つ。従って入射像は $S_3$ で次元5。
六行四列の最大小行列は期待余次元3を持ち、Eagon--Northcottにより商はCohen--Macaulayである。
階数3上は核線の消去で滑らかな入射と局所同型だから一般被約、従ってこの小行列スキームは被約である。
核に独立な三次式 $q,r$ がある場合、その最大公約式 $d$ に対し $qP_5+rP_5=dP_{8-\deg d}$
（最小公倍式で次元を数える）。$\deg d\le2$ より、これはちょうど $S_2$ の条件となる。
こちらも二次式に対する入射により次元3である。

射影の各ファイバーは中心を超平面に持つ線型空間内の、中心と交わらない射影スキームなので有限である。
像の次元は5で全射となり、Cohen--Macaulay性と標的の正則性から平坦である。
$h$ を三次式の空間の超平面、$H$ を入射の $Z$ 座標の超平面とすると、
$c(E)=(1+h)^{-6}$ により $\int H^5=\binom63=20$, $\int hH^4=\binom62=15$ となる。
階数3上は核線が逆射を与え、グラフはファイバーに一致する。
反対射影でも $q$ を固定した $p\mapsto pq$ の単射性から核は高々一次元で、同じく閉埋込みである。
階数2以下の点があれば、その上の $[q]$ が正次元となる。
\end{proof}
\begin{proof}[判別式と一般パラメータでの非正則Frobenius性]
$\pi(S_2)$ は次元3なので、一般直線はこれと交わらない。
入射上に四つの一般超平面を戻して滑らかな曲線を得る（標数零のBertini）。
入射の標準類は $-3H+2h$ なので、この曲線から直線への射の分岐の次数は
$\int (3H+2h)H^4=90$ である。一般直線上での有限平坦トレース判別式の位数は、
滑らかな曲線の分岐指数引く一の和であり、これで全体の次数を得る。
特に一般ファイバーは被約である。

三次交叉の次元4と四次交叉の次元1は、上述の隣接階数と一文字階数を固定すれば開条件である。
空でないことは単項式番号順に
\[
 T=(2,3,5,7,11,13,17,19,23,29,31),\qquad Y=(3,1,4,1,5,9,2,6,5)
\]
を使って有理数上で確認できる。
各向きで隣接二挿入の連結行列の階数は44、三次縮約を両側から挿入した連結行列の階数は31である。
全パラメータ空間は既約で、二条件の開集合は被約有限グラフの非空開集合と交わる。
定理\ref{thm:finite-intrinsic-stability-test}によりこれらは安定である。
計算は\path{computations/verify_binary_secant_obstructions.py}で再現できる。
\end{proof}
